\subsection{Convolutions}
\lecdate{(Midterm last thurs) Lec 10 - Feb 10 (Week 6)}
we say a fn is cptly supported if its supp is cpt.
\begin{defn}
	if $f,g:\bR\gto\bR$ cptly supp and riemann int'bl, then:
	\begin{equation*}
		(f*g)(x) \ = \ \int_{\bR}f(x-t)g(t)\d t
	\end{equation*}
	is well defn'd, and is called the \textbf{convolution} of $f$ and $g$.
\end{defn} \

\begin{defn}
	a seq $J_{n}:\bR\gto\bR$  of cptly supp riemann int'bl fn's are
	called an \textbf{approximate identity} if:
	\begin{enumerate}[(i)]
		\item $\displaystyle\int_{\bR}I_{n}(x)\d x=1$
		\item $\sup_{n}\displaystyle\int_{\bR}\abs{I_{n}(x)}\d x<\infty$
		\item $\displaystyle\lim_{n\sto\infty}\int_{\abs{x}>\delta}\abs{I_{n}(x)}\d x=0$
			for all $\delta>0$
	\end{enumerate}
\end{defn} \

\begin{thm}
	let $f:[a,b]\gto\bR$ be cts, $f(a)=f(b)=0$. define $\tilde{f}(x)$ as:
	\begin{equation*}
		\tilde{f}(x) \ = \
		\begin{cases}
			f(x) & x\in[a,b] \\ 0 & \trm{otw}
		\end{cases}
	\end{equation*}
	let $I_{n}$ be an approx id.
	then $f_{n}=\tilde{f}*I_{n}$ cvgs uniformly to $\tilde{f}$.
\end{thm} \

\begin{pf}[source=Primary Source Material]
	let $\ep>0$. easy to see: $\tilde{f}$ uni cts, so $\ex\delta>0$ s.t.
	\begin{equation*}
		\abs{\tilde{f}(x)-\tilde{f}(y)}<\ep \qquad \fa\abs{x-y}<\delta
	\end{equation*}
	then:
	\begin{align*}
		f_{n}(x)-\tilde{f}(x)
		& = \int\tilde{f}(x-y)I_{n}(y)\d y - \int\tilde{f}(x)I_{n}(y)\d y \\
		& = \int_{\bR}(\tilde{f}(x-y)-\tilde{f}(x))I_{n}(y)\d y \\
		& = \int_{\abs{y}<\delta}(\tilde{f}(x-y)-\tilde{f}(x))I_{n}(y)\d y
		+ \int_{\abs{y}>\delta}(\tilde{f}(x-y)-\tilde{f}(x))I_{n}(y)\d y
	\end{align*}
	first term:
	\begin{equation*}
		\leq \int_{\abs{y}<\delta}\abs{\tilde{f}(x-y)-\tilde{f}(x)}\abs{I_{n}(y)}\d y
		\leq \ep\cdot\sup_{n}\int_{\bR}\abs{I_{n}(y)}\d y
	\end{equation*}
	2nd term:
	\begin{equation*}
		\leq 2\norm{\tilde{f}}_{\infty}\int_{\abs{y}>\delta}\abs{I_{n}(y)}\d y
	\end{equation*}
	goes to 0 as $n\sto\infty$.
	thus, $\limsup_{n\sto\infty}\sup_{x\in\bR}\abs{f_{n}(x)-\tilde{f}(x)}\leq C\ep$.
\end{pf}

\newpage
\begin{thm}[title=Weierstrass Approximation Theorem]
	sps $f:[a,b]\gto\bR$ cts.
	then $\ex$ seq of polyn's $p_{n}$ s.t. $p_{n}\rightrightarrows f$ for $x\in[a,b]$.
\end{thm} \

\begin{pf}[source=Primary Source Material]
	by subtracting linear fn, wlog $f(a)=f(b)=0$.
	by rescaling, wlog $a=0$ and $b=1$.

	extend $f$ to $\bR$ with:
	\begin{equation*}
		\tilde{f}(x) \ = \
		\begin{cases}
			f(x) & x\in[0,1] \\ 0 & \trm{otw}
		\end{cases}
	\end{equation*}
	define $\beta_{n}(t):=b_{n}(1-t^{2})^{n}\bb{I}_{\abs{t}\leq1}$,
	where $\bb{I}$ indicator fn, and:
	\begin{equation*}
		b_{n} \ = \ \frac{1}{\int_{-1}^{1}(1-t^{2})^{n}\d t}
	\end{equation*}
	we claim the following:
	\begin{enumerate}[(1),itemsep=0pt]
		\item $\beta_{n}$ an approx id
		\item $f_{n}(x)=(\tilde{f}*\beta_{n})(x)$ a polyn $\fa n$ for $x\in[0,1]$
	\end{enumerate}
	together, these prove the statement.

	first, we prove claim 2. we have:
	\begin{equation*}
		f_{n}(x) \ = \ \int_{-1}^{1}\tilde{f}(x-t)\beta_{n}(t)\d t
	\end{equation*}
	by change of variables $u=x-t$, this is equal to:
	\begin{equation*}
		\int_{x-1}^{x+1}\tilde{f}(u)\beta_{n}(x-u)\d u
		= \int_{0}^{1}f(u)\beta_{n}(x-u)\d u
	\end{equation*}
	if $x\in[0,1]$ and $u\in[0,1]$, then $x-u\in[-1,1]$, so this is equal to:
	\begin{equation*}
		b_{n}\int_{0}^{1}f(u)(1-(x-u)^{2})^{n}\d u
	\end{equation*}
	this is a polyn in $x$, since this is equal to:
	\begin{equation*}
		\sum_{k,j}x^{j}\int_{0}^{1}f(u)u^{k}\d u
	\end{equation*}

	now, we prove claim 1. by defn:
	\begin{equation*}
		\int_{\bR}\beta_{n}(t)\d t
		\ = \ \int_{\bR}\abs{\beta_{n}(t)}\d t \ = \ 1
	\end{equation*}
	now, we show:
	\begin{equation*}
		\int_{\abs{t}>\delta}\beta_{n}(t)\d t \ \sto \ 0 \trm{ as } n\sto\infty
	\end{equation*}
	we find an upper bound for $b_{n}$:
	\begin{align*}
		\int_{-1}^{1}(1-t^{2})^{n}\d t
		& \geq \ \int_{-1/\sqrt{n}}^{1/\sqrt{n}}(1-t^{2})^{n}\d t \\
		& \geq \ \int_{-1/\sqrt{n}}^{1/\sqrt{n}}\left(1-\left(\frac{1}{\sqrt{n}}\right)^{2}\right)^{n}\d t \\
		& = \ \frac{2}{\sqrt{n}}\left(1-\frac{1}{n}\right)^{n} \\
		& \geq \ \frac{C}{\sqrt{n}}
	\end{align*}
	holds for all $n\geq2$. thus $b_{n}\leq C\sqrt{n}$, so:
	\begin{align*}
		\int_{\abs{t}>\delta}\beta_{n}(t)\d t
		& \ \leq \ C\sqrt{n}\int_{\abs{t}>\delta}(1-t^{2})^{n}\bb{I}_{\abs{t}\leq1} \\
		& \ \leq \ C\sqrt{n}(1-\delta^{2})^{n}
	\end{align*}
	this goes to 0 as $n\sto\infty$ as needed.

	we have proved both claims, thus the statement is proved.
\end{pf}

why are we defining algebras now lmfao

\begin{defn}
	we say $A\subseteq C_{0}(M)$ \textbf{separates pts} if $\fa p,q\in M$, $\ex f\in A$ s.t.
	$f(p)\neq f(q)$.

	we say $A$ \textbf{vanishes} at $x\in M$ if $f(x)=0$ $\fa f\in A$.
\end{defn} \

\begin{thm}[title=Stone-Weierstrass]
	if $M$ cpt, $A\subseteq C_{0}(M)$ an alg of fn's which separates pts and
	vanishes nowhere, then $\bar{A}=C_{0}(M)$.
\end{thm}



